\documentclass[11pt,a4paper]{article}
\usepackage[utf8]{inputenc}
\usepackage[french]{babel}
\usepackage{amsmath,amssymb,amsthm}
\usepackage{geometry}
\geometry{hmargin=2.5cm,vmargin=3cm}
\usepackage{tikz}
\usetikzlibrary{cd,positioning,shapes}
\usepackage{listings}
\usepackage{xcolor}

\lstset{
    language=Python,
    basicstyle=\ttfamily\small,
    keywordstyle=\color{blue}\bfseries,
    stringstyle=\color{red},
    commentstyle=\color{gray}\itshape,
    numbers=left,
    numberstyle=\tiny\color{gray},
    breaklines=true,
    frame=single,
    showstringspaces=false
}

\title{\Huge JALON 81 : TRANSFORMÉE DE FOURIER DANS $L^2$ ET THÉORÈME DE PLANCHEREL \\ \large Cours magistral, exercices corrigés et travaux pratiques}
\author{\Large Charles EDOU NZE}
\date{\today}

\begin{document}
\maketitle
\tableofcontents
\newpage

\part{Cours Magistral}
\section{Jalon 81 : Transformée de Fourier dans $L^2$ et Théorème de Plancherel}

\subsection*{Introduction}

La transformée de Fourier, initialement définie pour les fonctions de l'espace $L^1(\mathbb{R})$, est un outil puissant pour l'analyse harmonique. Cependant, l'espace $L^1$ ne possède pas de structure hilbertienne, contrairement à l'espace $L^2(\mathbb{R})$ des fonctions de carré intégrable, qui est fondamental en mécanique quantique et en théorie du signal (où le carré de la norme représente l'énergie).

Le défi majeur réside dans le fait que l'intégrale définissant la transformée de Fourier ne converge pas absolument pour une fonction $f \in L^2(\mathbb{R})$ quelconque. La construction de la transformée de Fourier sur $L^2(\mathbb{R})$ requiert donc une approche par prolongement continu. Ce passage est rendu possible par le théorème de Plancherel, qui établit que la transformée de Fourier, convenablement normalisée, agit comme une isométrie sur l'espace dense $L^1 \cap L^2$.

\subsection*{Théorème de Plancherel et Prolongement par Densité}

\subsubsection*{Définitions et Théorème de Plancherel}

Soit $\mathcal{S}(\mathbb{R})$ l'espace de Schwartz des fonctions à décroissance rapide. L'espace $\mathcal{S}(\mathbb{R})$ est dense à la fois dans $L^1(\mathbb{R})$ et dans $L^2(\mathbb{R})$. Pour $f \in L^1(\mathbb{R})$, la transformée de Fourier est définie par :
$$ \hat{f}(\xi) = \mathcal{F}(f)(\xi) = \int_{\mathbb{R}} f(t) e^{-i\xi t} dt $$

\textbf{Théorème (Plancherel-Parseval sur l'espace de Schwartz) :}
Pour toute fonction $f \in \mathcal{S}(\mathbb{R})$, sa transformée de Fourier $\hat{f}$ appartient également à $\mathcal{S}(\mathbb{R})$, et l'identité suivante (identité de Plancherel) est vérifiée :
$$ \|\hat{f}\|_{L^2}^2 = 2\pi \|f\|_{L^2}^2 $$
Autrement dit, $\int_{\mathbb{R}} |\hat{f}(\xi)|^2 d\xi = 2\pi \int_{\mathbb{R}} |f(t)|^2 dt$.

\textbf{Exemple 1 : Vérification pour la fonction gaussienne}
Considérons la fonction gaussienne normalisée $f(t) = e^{-t^2/2}$.
Sa norme $L^2$ au carré est :
$$ \|f\|_{L^2}^2 = \int_{\mathbb{R}} (e^{-t^2/2})^2 dt = \int_{\mathbb{R}} e^{-t^2} dt = \sqrt{\pi} $$
Sa transformée de Fourier est bien connue (caractère propre de la gaussienne) : $\hat{f}(\xi) = \sqrt{2\pi} e^{-\xi^2/2}$.
Calculons la norme $L^2$ au carré de $\hat{f}$ :
$$ \|\hat{f}\|_{L^2}^2 = \int_{\mathbb{R}} (\sqrt{2\pi} e^{-\xi^2/2})^2 d\xi = 2\pi \int_{\mathbb{R}} e^{-\xi^2} d\xi = 2\pi \sqrt{\pi} $$
On vérifie immédiatement que $\|\hat{f}\|_{L^2}^2 = 2\pi \|f\|_{L^2}^2$, conformément au théorème de Plancherel.

\subsubsection*{Le processus de prolongement}

L'application $\mathcal{F} : \mathcal{S}(\mathbb{R}) \to L^2(\mathbb{R})$ est donc une application linéaire continue vis-à-vis de la topologie de $L^2$ (à une constante $\sqrt{2\pi}$ près).

Puisque $\mathcal{S}(\mathbb{R})$ est un sous-espace dense de $L^2(\mathbb{R})$, le théorème de prolongement des applications uniformément continues permet d'étendre de manière unique cette application à l'espace complet $L^2(\mathbb{R})$.

\textbf{Théorème (Transformée de Fourier dans $L^2$) :}
Il existe un unique prolongement linéaire continu de la transformée de Fourier de $\mathcal{S}(\mathbb{R})$ (ou de $L^1(\mathbb{R}) \cap L^2(\mathbb{R})$) à $L^2(\mathbb{R})$, que nous noterons encore $\mathcal{F}$.
Pour tout $f \in L^2(\mathbb{R})$, $\mathcal{F}(f) \in L^2(\mathbb{R})$ et l'égalité $\|\mathcal{F}(f)\|_{L^2} = \sqrt{2\pi} \|f\|_{L^2}$ reste vraie.

Concrètement, pour $f \in L^2(\mathbb{R})$, la transformée de Fourier est définie comme la limite dans $L^2$ des transformées tronquées :
$$ \hat{f} = \lim_{R \to \infty}^{(L^2)} \int_{-R}^{R} f(t) e^{-i\xi t} dt $$
Cette limite est prise au sens de la norme $L^2$, c'est-à-dire que $\lim_{R \to \infty} \left\| \hat{f}(\cdot) - \int_{-R}^{R} f(t) e^{-i(\cdot)t} dt \right\|_{L^2} = 0$.

\textbf{Exemple 2 : Fonction porte et sinus cardinal}
Considérons la fonction indicatrice (ou "porte") $f(t) = \mathbf{1}_{[-a, a]}(t)$ pour $a > 0$. Clairement, $f \in L^1 \cap L^2$.
Sa transformée de Fourier se calcule directement :
$$ \hat{f}(\xi) = \int_{-a}^{a} e^{-i\xi t} dt = \left[ \frac{e^{-i\xi t}}{-i\xi} \right]_{-a}^{a} = \frac{e^{i\xi a} - e^{-i\xi a}}{i\xi} = \frac{2\sin(a\xi)}{\xi} = 2a \, \text{sinc}(a\xi) $$
Vérifions Plancherel.
L'énergie de $f$ est $\|f\|_{L^2}^2 = \int_{-a}^{a} 1^2 dt = 2a$.
L'énergie de $\hat{f}$ est $\|\hat{f}\|_{L^2}^2 = \int_{\mathbb{R}} \left( \frac{2\sin(a\xi)}{\xi} \right)^2 d\xi$.
Par le théorème de Plancherel, $\|\hat{f}\|_{L^2}^2 = 2\pi \times 2a = 4\pi a$.
Nous retrouvons ainsi élégamment l'intégrale de Dirichlet classique de manière indirecte : $\int_{\mathbb{R}} \frac{\sin^2(a\xi)}{\xi^2} d\xi = a\pi$.

\subsection*{Le Théorème d'Inversion et l'Isomorphisme}

Un résultat central de la théorie de Fourier sur $L^2$ est que l'opérateur de Fourier, à normalisation près, est un opérateur unitaire.

\textbf{Théorème (Inversion de Fourier dans $L^2$) :}
L'opérateur $\frac{1}{\sqrt{2\pi}}\mathcal{F}$ est un opérateur unitaire (bijectif et isométrique) de $L^2(\mathbb{R})$ sur lui-même.
En particulier, $\mathcal{F}$ est une bijection de $L^2(\mathbb{R})$ sur lui-même, et la formule d'inversion est vraie presque partout :
$$ f(t) = \frac{1}{2\pi} \int_{\mathbb{R}} \hat{f}(\xi) e^{i\xi t} d\xi $$
(L'intégrale étant comprise, là encore, au sens d'une limite $L^2$ de troncatures).

L'identité de Parseval généralisée pour le produit scalaire stipule que pour $f, g \in L^2(\mathbb{R})$ :
$$ \langle f, g \rangle_{L^2} = \frac{1}{2\pi} \langle \hat{f}, \hat{g} \rangle_{L^2} $$
soit
$$ \int_{\mathbb{R}} f(t) \overline{g(t)} dt = \frac{1}{2\pi} \int_{\mathbb{R}} \hat{f}(\xi) \overline{\hat{g}(\xi)} d\xi $$

\textbf{Exemple 3 : Produit scalaire de deux fonctions portes}
Soit $f(t) = \mathbf{1}_{[-1, 1]}(t)$ et $g(t) = \mathbf{1}_{[-2, 2]}(t)$.
Le produit scalaire direct dans $L^2$ est :
$$ \langle f, g \rangle = \int_{\mathbb{R}} \mathbf{1}_{[-1, 1]}(t) \mathbf{1}_{[-2, 2]}(t) dt = \int_{-1}^{1} 1 dt = 2 $$
En utilisant les transformées de Fourier, on a $\hat{f}(\xi) = \frac{2\sin(\xi)}{\xi}$ et $\hat{g}(\xi) = \frac{2\sin(2\xi)}{\xi}$.
L'identité de Parseval nous donne :
$$ \frac{1}{2\pi} \int_{\mathbb{R}} \frac{2\sin(\xi)}{\xi} \frac{2\sin(2\xi)}{\xi} d\xi = 2 $$
Ce qui permet d'évaluer une intégrale non triviale : $\int_{\mathbb{R}} \frac{\sin(\xi)\sin(2\xi)}{\xi^2} d\xi = \pi$.


\begin{center}
\begin{tikzpicture}[scale=1.5]
    % Domain Temporel
    \draw[thick, ->] (-2,0) -- (2,0) node[right] {$t$};
    \draw[thick, ->] (0,-0.5) -- (0,2) node[above] {$f(t)$};
    \draw[blue, thick, domain=-1.5:1.5, samples=100] plot (\x, {exp(-abs(\x))});

    % Arrow
    \draw[thick, ->] (2.5, 0.5) -- (3.5, 0.5) node[midway, above] {$\mathcal{F}$};

    % Domain Frequentiel
    \begin{scope}[shift={(6,0)}]
        \draw[thick, ->] (-2,0) -- (2,0) node[right] {$\xi$};
        \draw[thick, ->] (0,-0.5) -- (0,2.5) node[above] {$\hat{f}(\xi)$};
        \draw[red, thick, domain=-2:2, samples=100] plot (\x, {2/(1+\x*\x)});
    \end{scope}

    \node at (2, -1) {Énergie Temporelle: $\int |f(t)|^2 dt$};
    \node at (6, -1) {Énergie Fréquentielle: $\frac{1}{2\pi} \int |\hat{f}(\xi)|^2 d\xi$};
    \node at (4, -1.5) {Théorème de Plancherel: Égalité des Énergies};
\end{tikzpicture}
\end{center}

\subsection*{Démonstrations Pas-à-Pas}

\subsubsection*{Preuve du théorème de Plancherel pour la classe de Schwartz}

Soient $f, g \in \mathcal{S}(\mathbb{R})$. Nous allons montrer l'identité de Parseval $\langle f, g \rangle = \frac{1}{2\pi} \langle \hat{f}, \hat{g} \rangle$. En prenant $f=g$, nous obtiendrons le théorème de Plancherel.

1. \textbf{Étape 1 : Formule de Dualité}
Définissons la fonction d'interversion. Pour $f, g \in \mathcal{S}(\mathbb{R})$, considérons l'intégrale double :
$$ I = \int_{\mathbb{R}} \int_{\mathbb{R}} f(x) g(y) e^{-ixy} dx dy $$
Comme $f$ et $g$ sont à décroissance rapide, la fonction intégrande $(x,y) \mapsto f(x)g(y)e^{-ixy}$ est dominée par $|f(x)||g(y)|$, qui est intégrable sur $\mathbb{R}^2$.
Le théorème de Fubini s'applique, et nous pouvons intégrer dans n'importe quel ordre.
Intégration d'abord par rapport à $x$ :
$$ I = \int_{\mathbb{R}} g(y) \left( \int_{\mathbb{R}} f(x) e^{-ixy} dx \right) dy = \int_{\mathbb{R}} g(y) \hat{f}(y) dy $$
Intégration d'abord par rapport à $y$ :
$$ I = \int_{\mathbb{R}} f(x) \left( \int_{\mathbb{R}} g(y) e^{-ixy} dy \right) dx = \int_{\mathbb{R}} f(x) \hat{g}(x) dx $$
Nous obtenons ainsi la formule d'échange :
$$ \int_{\mathbb{R}} \hat{f}(y) g(y) dy = \int_{\mathbb{R}} f(x) \hat{g}(x) dx $$

2. \textbf{Étape 2 : Application à la fonction de Gauss}
Soit $f \in \mathcal{S}(\mathbb{R})$. Soit $h(x) = e^{-a x^2 / 2}$ (avec $a > 0$).
Sa transformée est $\hat{h}(\xi) = \sqrt{\frac{2\pi}{a}} e^{-\xi^2 / (2a)}$.
Appliquons la formule d'échange avec $f$ et $h$ :
$$ \int_{\mathbb{R}} \hat{f}(y) e^{-a y^2 / 2} dy = \int_{\mathbb{R}} f(x) \sqrt{\frac{2\pi}{a}} e^{-x^2 / (2a)} dx $$

3. \textbf{Étape 3 : Passage à la limite pour construire l'inversion}
Multiplions les deux côtés par $\frac{1}{\sqrt{2\pi}}$ (ou considérons cela en posant $x = \sqrt{a} u$ dans le second membre).
$$ \int_{\mathbb{R}} \hat{f}(y) e^{-a y^2 / 2} dy = \sqrt{2\pi} \int_{\mathbb{R}} f(\sqrt{a}u) e^{-u^2/2} du $$
Prenons la limite quand $a \to 0^+$.
À gauche, $\lim_{a \to 0} e^{-a y^2 / 2} = 1$. Comme $\hat{f} \in \mathcal{S} \subset L^1$, le théorème de convergence dominée s'applique.
À droite, $f(\sqrt{a}u) \to f(0)$. Comme $f$ est bornée (étant dans $\mathcal{S}$), la fonction intégrande est dominée par $M e^{-u^2/2}$, intégrable.
On obtient :
$$ \int_{\mathbb{R}} \hat{f}(y) dy = \sqrt{2\pi} f(0) \int_{\mathbb{R}} e^{-u^2/2} du = \sqrt{2\pi} f(0) \times \sqrt{2\pi} = 2\pi f(0) $$
On en déduit que $f(0) = \frac{1}{2\pi} \int_{\mathbb{R}} \hat{f}(y) dy$.
Par translation, en appliquant ce résultat à $f_x(t) = f(t+x)$ dont la transformée est $\hat{f}_x(\xi) = \hat{f}(\xi)e^{i\xi x}$, on obtient la formule d'inversion :
$$ f(x) = \frac{1}{2\pi} \int_{\mathbb{R}} \hat{f}(\xi) e^{i\xi x} d\xi $$

4. \textbf{Étape 4 : Identité de Parseval et Plancherel}
Nous avons, d'après la formule d'inversion, pour toute fonction $h \in \mathcal{S}$ :
$$ h(x) = \frac{1}{2\pi} \int_{\mathbb{R}} \hat{h}(\xi) e^{i\xi x} d\xi $$
Soient $f, g \in \mathcal{S}$. Posons $h = f * g^*$ où $g^*(x) = \overline{g(-x)}$.
On sait que la transformée d'une convolution est le produit des transformées : $\hat{h} = \hat{f} \widehat{g^*}$.
Et $\widehat{g^*}(\xi) = \overline{\hat{g}(\xi)}$.
Ainsi, $\hat{h}(\xi) = \hat{f}(\xi) \overline{\hat{g}(\xi)}$.
Évaluons $h$ en 0 de deux manières.
D'une part, par définition de la convolution :
$$ h(0) = (f * g^*)(0) = \int_{\mathbb{R}} f(t) g^*(0-t) dt = \int_{\mathbb{R}} f(t) \overline{g(t)} dt = \langle f, g \rangle $$
D'autre part, par la formule d'inversion pour $h$ en $x=0$ :
$$ h(0) = \frac{1}{2\pi} \int_{\mathbb{R}} \hat{h}(\xi) e^{i\xi \cdot 0} d\xi = \frac{1}{2\pi} \int_{\mathbb{R}} \hat{f}(\xi) \overline{\hat{g}(\xi)} d\xi = \frac{1}{2\pi} \langle \hat{f}, \hat{g} \rangle $$
En égalant les deux expressions, nous obtenons l'identité de Parseval. En prenant $f = g$, nous obtenons l'identité de Plancherel $\|\hat{f}\|_{L^2}^2 = 2\pi \|f\|_{L^2}^2$, ce qui achève la démonstration.

\subsection*{Applications}

\subsubsection*{En Physique : Densité Spectrale de Puissance}

L'identité de Plancherel est la pierre angulaire du traitement du signal et de la physique quantique. Pour un signal temporel $s(t)$ modélisant par exemple un courant électrique traversant une résistance de $1\,\Omega$, la puissance instantanée dissipée est proportionnelle à $|s(t)|^2$.
L'énergie totale dissipée sur tout le temps est $E = \int_{\mathbb{R}} |s(t)|^2 dt$.
Le théorème de Plancherel énonce que $E = \frac{1}{2\pi} \int_{\mathbb{R}} |\hat{s}(\xi)|^2 d\xi$.
La quantité $|\hat{s}(\xi)|^2$ est appelée la \textbf{densité spectrale d'énergie}. Elle quantifie la distribution de l'énergie du signal à travers ses différentes fréquences composantes $\xi$. Cette équivalence stipule le principe physique fondamental de conservation de l'énergie entre le domaine temporel et le domaine fréquentiel.

\subsubsection*{En Intelligence Artificielle et Analyse de Données}

Dans le domaine du Machine Learning et des réseaux de neurones, la transformée de Fourier dans $L^2$ intervient dans :
1.  \textbf{L'analyse des opérateurs de convolution :} Dans les réseaux CNN, la convolution par un filtre $w$ dans l'espace temporel équivaut à une multiplication point par point par $\hat{w}$ dans l'espace de Fourier. Plancherel assure que la norme spectrale de l'opérateur (déterminant la constante de Lipschitz du réseau, cruciale pour la robustesse et les WGANs) est directement reliée au supremum de $|\hat{w}(\xi)|$.
2.  \textbf{Traitement Audio et Parole :} Les algorithmes d'extraction de caractéristiques comme le MFCC (Mel-Frequency Cepstral Coefficients) s'appuient sur le calcul de la densité spectrale d'énergie (le module au carré de la transformée de Fourier, validé par Plancherel) pour capturer les "formants" du signal vocal, avant d'appliquer des échelles logarithmiques imitant la perception humaine.


\newpage
\part{Exercices d'Application et de Concours}
\subsection*{Exercice 1 : Calcul d'énergie via Plancherel \quad $\bigstar\star\star\star\star$}

\textbf{Énoncé :}
Considérons le signal $f(t) = e^{-a|t|}$ avec $a > 0$.
1. Calculer l'énergie totale du signal dans le domaine temporel : $E = \|f\|_{L^2}^2 = \int_{-\infty}^{+\infty} |f(t)|^2 dt$.
2. Calculer la transformée de Fourier $\hat{f}(\xi)$.
3. Calculer l'énergie totale dans le domaine fréquentiel : $\frac{1}{2\pi} \|\hat{f}\|_{L^2}^2 = \frac{1}{2\pi} \int_{-\infty}^{+\infty} |\hat{f}(\xi)|^2 d\xi$.
4. Vérifier l'égalité de Plancherel-Parseval.

\textbf{Correction :}
1. \textbf{Énergie temporelle :}
La fonction $f(t) = e^{-a|t|}$ est paire. Son module au carré est $|f(t)|^2 = e^{-2a|t|}$.
$$ E = \int_{-\infty}^{+\infty} e^{-2a|t|} dt = 2 \int_{0}^{+\infty} e^{-2at} dt $$
En primitivant :
$$ E = 2 \left[ \frac{e^{-2at}}{-2a} \right]_0^{+\infty} = 2 \left( 0 - \frac{1}{-2a} \right) = \frac{1}{a} $$

2. \textbf{Transformée de Fourier :}
Par définition, $\hat{f}(\xi) = \int_{-\infty}^{+\infty} e^{-a|t|} e^{-i\xi t} dt$.
En séparant l'intégrale :
$$ \hat{f}(\xi) = \int_{-\infty}^{0} e^{at} e^{-i\xi t} dt + \int_{0}^{+\infty} e^{-at} e^{-i\xi t} dt $$
$$ \hat{f}(\xi) = \int_{-\infty}^{0} e^{(a-i\xi)t} dt + \int_{0}^{+\infty} e^{-(a+i\xi)t} dt $$
$$ \hat{f}(\xi) = \left[ \frac{e^{(a-i\xi)t}}{a-i\xi} \right]_{-\infty}^0 + \left[ \frac{e^{-(a+i\xi)t}}{-(a+i\xi)} \right]_0^{+\infty} $$
Comme $a>0$, $\lim_{t \to -\infty} e^{at} = 0$ et $\lim_{t \to +\infty} e^{-at} = 0$.
$$ \hat{f}(\xi) = \frac{1}{a-i\xi} - 0 + 0 - \frac{1}{-(a+i\xi)} = \frac{1}{a-i\xi} + \frac{1}{a+i\xi} $$
En réduisant au même dénominateur :
$$ \hat{f}(\xi) = \frac{a+i\xi + a-i\xi}{(a-i\xi)(a+i\xi)} = \frac{2a}{a^2 + \xi^2} $$

3. \textbf{Énergie fréquentielle :}
Calculons l'intégrale du module au carré :
$$ I = \int_{-\infty}^{+\infty} \left( \frac{2a}{a^2 + \xi^2} \right)^2 d\xi = 4a^2 \int_{-\infty}^{+\infty} \frac{1}{(a^2 + \xi^2)^2} d\xi $$
Pour évaluer cette intégrale, posons le changement de variable $\xi = a \tan \theta$, avec $d\xi = a(1+\tan^2\theta) d\theta = \frac{a}{\cos^2\theta} d\theta$.
Les bornes $-\infty$ et $+\infty$ deviennent $-\frac{\pi}{2}$ et $\frac{\pi}{2}$.
De plus, $a^2 + \xi^2 = a^2(1+\tan^2\theta) = \frac{a^2}{\cos^2\theta}$.
Donc $\frac{1}{(a^2+\xi^2)^2} = \frac{\cos^4\theta}{a^4}$.
$$ I = 4a^2 \int_{-\pi/2}^{\pi/2} \frac{\cos^4\theta}{a^4} \frac{a}{\cos^2\theta} d\theta = \frac{4a^3}{a^4} \int_{-\pi/2}^{\pi/2} \cos^2\theta d\theta $$
$$ I = \frac{4}{a} \int_{-\pi/2}^{\pi/2} \frac{1+\cos(2\theta)}{2} d\theta = \frac{2}{a} \left[ \theta + \frac{\sin(2\theta)}{2} \right]_{-\pi/2}^{\pi/2} $$
$$ I = \frac{2}{a} \left( \left(\frac{\pi}{2} + 0\right) - \left(-\frac{\pi}{2} + 0\right) \right) = \frac{2}{a} \times \pi = \frac{2\pi}{a} $$
L'énergie fréquentielle normalisée est donc $\frac{1}{2\pi} I = \frac{1}{2\pi} \times \frac{2\pi}{a} = \frac{1}{a}$.

4. \textbf{Vérification :}
On constate bien que $\|f\|_{L^2}^2 = \frac{1}{a} = \frac{1}{2\pi} \|\hat{f}\|_{L^2}^2$, validant ainsi l'identité de Plancherel-Parseval pour ce signal.



\subsection*{Exercice 2 : Théorème d'inversion pour des fonctions affines par morceaux \quad $\bigstar\bigstar\star\star\star$}

\textbf{Énoncé :}
Soit la fonction "triangle" $f(t) = \max(1 - |t|, 0)$.
1. Montrer que $f \in L^1(\mathbb{R}) \cap L^2(\mathbb{R})$.
2. Calculer sa transformée de Fourier $\hat{f}(\xi)$.
3. Utiliser le théorème d'inversion de Fourier pour calculer l'intégrale $\int_{-\infty}^{+\infty} \frac{\sin^2(\xi/2)}{(\xi/2)^2} d\xi$.

\textbf{Correction :}
1. \textbf{Appartenance à $L^1$ et $L^2$ :}
La fonction est continue, à support compact $[-1, 1]$, et bornée par $1$.
Elle est trivialement dans $L^1(\mathbb{R})$ : $\|f\|_1 = \int_{-1}^1 (1-|t|) dt = 2 \int_0^1 (1-t) dt = 2 \left[ t - \frac{t^2}{2} \right]_0^1 = 1$.
Elle est aussi dans $L^2(\mathbb{R})$ : $\|f\|_2^2 = 2 \int_0^1 (1-t)^2 dt = 2 \left[ \frac{-(1-t)^3}{3} \right]_0^1 = \frac{2}{3} < \infty$.

2. \textbf{Transformée de Fourier :}
$f$ peut s'écrire comme l'auto-convolution de la fonction porte $p(t) = \mathbf{1}_{[-1/2, 1/2]}(t)$.
En effet, $(p * p)(t) = \int p(\tau) p(t-\tau) d\tau = \int_{-1/2}^{1/2} \mathbf{1}_{[-1/2, 1/2]}(t-\tau) d\tau$.
L'intégrande est non nul si $-1/2 \leq t-\tau \leq 1/2$, soit $t-1/2 \leq \tau \leq t+1/2$.
L'intersection de $[-1/2, 1/2]$ et $[t-1/2, t+1/2]$ a pour longueur $\max(1-|t|, 0)$. Donc $f = p * p$.
Par propriété de la transformée de Fourier, $\mathcal{F}(p * p)(\xi) = (\mathcal{F}(p)(\xi))^2$.
Or, $\mathcal{F}(p)(\xi) = \int_{-1/2}^{1/2} e^{-i\xi t} dt = \left[ \frac{e^{-i\xi t}}{-i\xi} \right]_{-1/2}^{1/2} = \frac{e^{i\xi/2} - e^{-i\xi/2}}{i\xi} = \frac{2\sin(\xi/2)}{\xi}$.
Donc, $\hat{f}(\xi) = \left( \frac{2\sin(\xi/2)}{\xi} \right)^2 = \left( \frac{\sin(\xi/2)}{\xi/2} \right)^2 = \text{sinc}^2(\xi/2)$.

3. \textbf{Inversion de Fourier :}
Comme $\hat{f}(\xi) \ge 0$ et $\int \hat{f}(\xi) d\xi$ est lié à $f(0)$, vérifions l'intégrabilité de $\hat{f}$.
$\hat{f}(\xi) \sim 1$ en $0$, et $\hat{f}(\xi) \le \frac{4}{\xi^2}$ à l'infini, donc $\hat{f} \in L^1(\mathbb{R})$.
On peut donc appliquer le théorème d'inversion ponctuelle pour toute fonction $f$ continue dont la transformée est dans $L^1$ :
$$ f(t) = \frac{1}{2\pi} \int_{-\infty}^{+\infty} \hat{f}(\xi) e^{i\xi t} d\xi $$
Évaluons cette identité en $t=0$ :
$$ f(0) = \frac{1}{2\pi} \int_{-\infty}^{+\infty} \hat{f}(\xi) e^0 d\xi = \frac{1}{2\pi} \int_{-\infty}^{+\infty} \frac{\sin^2(\xi/2)}{(\xi/2)^2} d\xi $$
Or $f(0) = \max(1-0, 0) = 1$.
On en déduit donc immédiatement :
$$ \int_{-\infty}^{+\infty} \frac{\sin^2(\xi/2)}{(\xi/2)^2} d\xi = 2\pi $$



\subsection*{Exercice 3 : Continuité et isométrie de l'opérateur de Fourier \quad $\bigstar\bigstar\star\star\star$}

\textbf{Énoncé :}
Montrer que si l'on munit $L^1 \cap L^2$ de la norme $\|\cdot\|_{L^2}$, l'opérateur $\mathcal{F} : (L^1 \cap L^2, \|\cdot\|_2) \to (L^2, \|\cdot\|_2)$ est uniformément continu.

\textbf{Correction :}
Soit l'opérateur linéaire $\mathcal{F}$ restreint à $L^1 \cap L^2$, sous-espace de $L^2$.
Par l'identité de Plancherel pour les fonctions de l'espace de Schwartz $\mathcal{S}$ (qui est dense dans $L^1 \cap L^2$ pour la norme $\|\cdot\|_2$), nous savons que :
$$ \forall \phi \in \mathcal{S}, \quad \|\mathcal{F}(\phi)\|_2 = \sqrt{2\pi} \|\phi\|_2 $$
Soit $f \in L^1 \cap L^2$. Par densité, il existe une suite $(\phi_n)$ dans $\mathcal{S}$ telle que $\|\phi_n - f\|_2 \to 0$ et $\|\phi_n - f\|_1 \to 0$.
Par la continuité de $\mathcal{F}$ sur $L^1$ vers $L^\infty$, $\mathcal{F}(\phi_n) \to \mathcal{F}(f)$ ponctuellement (et uniformément).
Le lemme de Fatou permet de passer à la limite dans l'identité hilbertienne :
$$ \|\mathcal{F}(f)\|_2^2 \le \liminf_{n \to \infty} \|\mathcal{F}(\phi_n)\|_2^2 = \liminf_{n \to \infty} 2\pi \|\phi_n\|_2^2 = 2\pi \|f\|_2^2 $$
En fait, l'égalité stricte est maintenue (ce qui peut être prouvé en utilisant la complétude et la continuité du produit scalaire).
Dès lors, pour toute fonction $f \in L^1 \cap L^2$, l'opérateur satisfait :
$$ \|\mathcal{F}(f)\|_2 \le \sqrt{2\pi} \|f\|_2 $$
Pour la continuité uniforme, prenons $f, g \in L^1 \cap L^2$.
La linéarité de l'opérateur permet d'écrire :
$$ \|\mathcal{F}(f) - \mathcal{F}(g)\|_2 = \|\mathcal{F}(f - g)\|_2 $$
En appliquant la majoration de la norme à la fonction $h = f - g \in L^1 \cap L^2$ :
$$ \|\mathcal{F}(f) - \mathcal{F}(g)\|_2 \le \sqrt{2\pi} \|f - g\|_2 $$
Cette inégalité lipschitzienne montre que l'opérateur est globalement lipschitzien de rapport $\sqrt{2\pi}$, donc a fortiori uniformément continu. Cela garantit l'existence d'un unique prolongement continu à l'espace complet $L^2$.



\subsection*{Exercice 4 : Produit de convolution et Parseval \quad $\bigstar\bigstar\bigstar\star\star$}

\textbf{Énoncé :}
1. Montrer que si $f, g \in L^2(\mathbb{R})$, alors $f * g$ est une fonction continue et bornée.
2. Évaluer $\int_{\mathbb{R}} \frac{\sin(a\xi)\sin(b\xi)}{\xi^2} d\xi$ pour $a > 0$ et $b > 0$.

\textbf{Correction :}
1. \textbf{Continuité de la convolution :}
Pour $f, g \in L^2$, on considère $(f * g)(x) = \int_{\mathbb{R}} f(t) g(x-t) dt$.
Soit $g_x(t) = g(x-t)$. Par l'inégalité de Cauchy-Schwarz :
$$ |(f * g)(x)| \le \int_{\mathbb{R}} |f(t)| |g(x-t)| dt \le \|f\|_2 \|g_x\|_2 $$
L'invariance par translation de la mesure de Lebesgue implique $\|g_x\|_2 = \|g\|_2$.
D'où $|(f * g)(x)| \le \|f\|_2 \|g\|_2$ pour tout $x$. La fonction est donc uniformément bornée.
Pour la continuité, considérons la différence :
$$ |(f * g)(x+h) - (f * g)(x)| = \left| \int_{\mathbb{R}} f(t) [g(x+h-t) - g(x-t)] dt \right| $$
Par Cauchy-Schwarz :
$$ |(f * g)(x+h) - (f * g)(x)| \le \|f\|_2 \|g_{x+h} - g_x\|_2 = \|f\|_2 \|g_{-h} - g\|_2 $$
La continuité en moyenne d'ordre $p$ (pour $p=2$) stipule que $\lim_{h \to 0} \|g_{-h} - g\|_2 = 0$.
Par conséquent, $\lim_{h \to 0} (f * g)(x+h) = (f * g)(x)$, ce qui démontre la continuité uniforme de $f * g$.

2. \textbf{Calcul d'intégrale via Parseval :}
Posons $f(t) = \mathbf{1}_{[-a, a]}(t)$ et $g(t) = \mathbf{1}_{[-b, b]}(t)$. Les deux fonctions sont dans $L^1 \cap L^2$.
Leurs transformées de Fourier sont $\hat{f}(\xi) = \frac{2\sin(a\xi)}{\xi}$ et $\hat{g}(\xi) = \frac{2\sin(b\xi)}{\xi}$.
L'identité de Parseval s'écrit :
$$ \langle \hat{f}, \hat{g} \rangle_{L^2} = 2\pi \langle f, g \rangle_{L^2} $$
$$ \int_{\mathbb{R}} \frac{4\sin(a\xi)\sin(b\xi)}{\xi^2} d\xi = 2\pi \int_{\mathbb{R}} \mathbf{1}_{[-a, a]}(t) \mathbf{1}_{[-b, b]}(t) dt $$
Le produit des fonctions indicatrices est l'indicatrice de l'intersection des intervalles $[-a, a] \cap [-b, b] = [-\min(a,b), \min(a,b)]$.
$$ \int_{\mathbb{R}} \frac{4\sin(a\xi)\sin(b\xi)}{\xi^2} d\xi = 2\pi \int_{-\min(a,b)}^{\min(a,b)} 1 dt = 2\pi (2\min(a,b)) = 4\pi\min(a,b) $$
En simplifiant par 4 :
$$ \int_{\mathbb{R}} \frac{\sin(a\xi)\sin(b\xi)}{\xi^2} d\xi = \pi \min(a,b) $$



\subsection*{Exercice 5 : Dualité temporelle-fréquentielle et inégalité d'Heisenberg \quad $\bigstar\bigstar\bigstar\star\star$}

\textbf{Énoncé :}
Considérons l'espace de Hilbert $L^2(\mathbb{R})$. Soit $f \in \mathcal{S}(\mathbb{R})$ telle que $\|f\|_2 = 1$.
1. Rappeler les expressions de la variance en position $V_x$ et la variance en impulsion (fréquence) $V_\xi$, pour une fonction centrée en $0$ dans les deux domaines.
2. Démontrer le principe d'incertitude d'Heisenberg : $V_x V_\xi \ge \frac{1}{4}$.
\textit{(Indication : Utiliser l'intégration par parties sur l'intégrale $\int t f(t) f'(t) dt$ et Cauchy-Schwarz)}

\textbf{Correction :}
1. \textbf{Variances :}
Comme $\|f\|_2 = 1$ et l'identité de Plancherel stipule $\|\hat{f}\|_2^2 = 2\pi$, la densité de probabilité dans l'espace des positions est $|f(t)|^2$, et dans l'espace des fréquences $\frac{1}{2\pi} |\hat{f}(\xi)|^2$.
Pour des densités centrées en zéro :
$$ V_x = \int_{\mathbb{R}} t^2 |f(t)|^2 dt = \| t \mapsto t f(t) \|_2^2 $$
$$ V_\xi = \frac{1}{2\pi} \int_{\mathbb{R}} \xi^2 |\hat{f}(\xi)|^2 d\xi $$
Or, on sait que $\mathcal{F}(f')(\xi) = i\xi \hat{f}(\xi)$. Par suite, $|\mathcal{F}(f')(\xi)|^2 = \xi^2 |\hat{f}(\xi)|^2$.
En appliquant Plancherel à $f'$ : $\frac{1}{2\pi} \int |\mathcal{F}(f')(\xi)|^2 d\xi = \int |f'(t)|^2 dt$.
Donc, $V_\xi = \|f'\|_2^2 = \int_{\mathbb{R}} |f'(t)|^2 dt$.

2. \textbf{Démonstration de l'inégalité d'Heisenberg :}
Considérons l'intégrale de $t (f(t) \overline{f'(t)} + \overline{f(t)} f'(t)) = t \frac{d}{dt} |f(t)|^2$.
Intégrons par parties sur $[-R, R]$ :
$$ \int_{-R}^{R} t \frac{d}{dt} |f(t)|^2 dt = \left[ t |f(t)|^2 \right]_{-R}^{R} - \int_{-R}^{R} |f(t)|^2 dt $$
Comme $f \in \mathcal{S}$, $|f(t)|^2$ décroît plus vite que tout polynôme, donc $\lim_{R \to \infty} R|f(\pm R)|^2 = 0$.
Par passage à la limite :
$$ \int_{\mathbb{R}} t (f(t) \overline{f'(t)} + \overline{f(t)} f'(t)) dt = - \int_{\mathbb{R}} |f(t)|^2 dt = - \|f\|_2^2 = -1 $$
La quantité sous l'intégrale de gauche est $2 \text{Re}(t f(t) \overline{f'(t)})$.
Ainsi, $2 \text{Re} \int_{\mathbb{R}} t f(t) \overline{f'(t)} dt = -1$.
Le module d'un nombre réel est égal au module de sa partie réelle (au signe près), et la partie réelle est inférieure au module du nombre complexe global :
$$ 1 = \left| 2 \text{Re} \int_{\mathbb{R}} t f(t) \overline{f'(t)} dt \right| \le 2 \int_{\mathbb{R}} |t f(t) \overline{f'(t)}| dt $$
Par l'inégalité de Cauchy-Schwarz appliquée aux fonctions $t f(t)$ et $f'(t)$ :
$$ 1 \le 2 \left( \int_{\mathbb{R}} |t f(t)|^2 dt \right)^{1/2} \left( \int_{\mathbb{R}} |f'(t)|^2 dt \right)^{1/2} = 2 \sqrt{V_x} \sqrt{V_\xi} $$
En élevant au carré et en divisant par 4, nous obtenons l'inégalité stricte :
$$ V_x V_\xi \ge \frac{1}{4} $$
Cela signifie physiquement qu'il est impossible qu'un signal soit arbitrairement concentré simultanément dans le temps (petite $V_x$) et dans la fréquence (petite $V_\xi$).



\subsection*{Exercice 6 : Résolution d'une équation différentielle via Fourier \quad $\bigstar\bigstar\bigstar\star\star$}

\textbf{Énoncé :}
Résoudre l'équation différentielle suivante dans $L^2(\mathbb{R})$ :
$$ -u''(t) + u(t) = e^{-|t|} $$

\textbf{Correction :}
Supposons que $u \in L^2$ est solution et que ses dérivées ont un sens tel que l'on puisse appliquer la transformée de Fourier.
Appliquons la transformée de Fourier (dans $L^2$, par prolongement de $L^1$) des deux côtés de l'équation.
Notons $\hat{u} = \mathcal{F}(u)$.
On a $\mathcal{F}(u'')(\xi) = (i\xi)^2 \hat{u}(\xi) = -\xi^2 \hat{u}(\xi)$.
L'équation devient dans le domaine fréquentiel :
$$ -(-\xi^2 \hat{u}(\xi)) + \hat{u}(\xi) = \mathcal{F}(e^{-|t|})(\xi) $$
$$ (\xi^2 + 1)\hat{u}(\xi) = \mathcal{F}(e^{-|t|})(\xi) $$
Nous avons calculé précédemment que pour $a=1$, $\mathcal{F}(e^{-|t|})(\xi) = \frac{2}{1+\xi^2}$.
L'équation algébrique est donc :
$$ (\xi^2 + 1)\hat{u}(\xi) = \frac{2}{1+\xi^2} $$
D'où l'on isole $\hat{u}$ :
$$ \hat{u}(\xi) = \frac{2}{(1+\xi^2)^2} $$
Pour revenir dans le domaine temporel, on remarque que $\frac{2}{(1+\xi^2)^2}$ peut être vu via la propriété de dérivation par rapport aux fréquences, ou par la convolution.
En effet, $\hat{u}(\xi) = \frac{1}{2} \times \frac{2}{1+\xi^2} \times \frac{2}{1+\xi^2} = \frac{1}{2} \mathcal{F}(e^{-|t|})(\xi) \mathcal{F}(e^{-|t|})(\xi)$.
Par la propriété liant produit fréquentiel et convolution temporelle :
$$ \mathcal{F}(u) = \frac{1}{2} \mathcal{F}(e^{-|\cdot|} * e^{-|\cdot|}) $$
L'inversion donne $u(t) = \frac{1}{2} (e^{-|\cdot|} * e^{-|\cdot|})(t)$.
Calculons cette convolution :
$$ v(t) = \int_{-\infty}^{+\infty} e^{-|\tau|} e^{-|t-\tau|} d\tau $$
Pour $t \ge 0$, séparons l'intégrale en trois régions par rapport aux valeurs absolues :
\textbullet\ Pour $\tau < 0$, $v_1(t) = \int_{-\infty}^0 e^\tau e^{-(t-\tau)} d\tau = e^{-t} \int_{-\infty}^0 e^{2\tau} d\tau = e^{-t} \left[\frac{e^{2\tau}}{2}\right]_{-\infty}^0 = \frac{e^{-t}}{2}$. \\
\textbullet\ Pour $0 \le \tau \le t$, $v_2(t) = \int_0^t e^{-\tau} e^{-(t-\tau)} d\tau = \int_0^t e^{-t} d\tau = t e^{-t}$. \\
\textbullet\ Pour $\tau > t$, $v_3(t) = \int_t^{+\infty} e^{-\tau} e^{t-\tau} d\tau = e^t \int_t^{+\infty} e^{-2\tau} d\tau = e^t \left[ \frac{e^{-2\tau}}{-2} \right]_t^{+\infty} = e^t \frac{e^{-2t}}{2} = \frac{e^{-t}}{2}$. \\
Sommons les trois contributions pour $t \ge 0$ :
$$ v(t) = \frac{e^{-t}}{2} + t e^{-t} + \frac{e^{-t}}{2} = (1+t)e^{-t} $$
Par symétrie évidente de la convolution de fonctions paires, $v(-t) = v(t)$.
Donc $v(t) = (1+|t|)e^{-|t|}$.
Finalement, $u(t) = \frac{1}{2} v(t) = \frac{1}{2} (1+|t|)e^{-|t|}$.
Cette fonction est bien dans $L^2(\mathbb{R})$, valide notre résolution et prouve la surjectivité de l'opérateur pour ce membre de droite.



\subsection*{Exercice 7 : Calcul d'une intégrale semi-convergente par l'isométrie \quad $\bigstar\bigstar\bigstar\bigstar\star$}

\textbf{Énoncé :}
Montrer, en utilisant l'identité de Plancherel, la valeur de l'intégrale $\int_0^{+\infty} \left(\frac{\sin(x^2)}{x}\right)^2 dx$.

\textbf{Correction :}
L'intégrale posée, avec un carré dans l'argument du sinus, ne correspond pas directement à la transformée d'une fonction porte basique. Posons $y = x^2$, alors $dx = \frac{dy}{2\sqrt{y}}$.
L'intégrale devient $J = \int_0^{+\infty} \frac{\sin^2(y)}{y} \frac{dy}{2\sqrt{y}} = \frac{1}{2} \int_0^{+\infty} y^{-3/2} \sin^2(y) dy$. Cela ne simplifie pas l'usage de Parseval.
Utilisons plutôt la transformée de Fourier de la fonction triangle.
Nous savons que pour $f(t) = \max(1-|t|, 0)$, $\hat{f}(\xi) = \frac{\sin^2(\xi/2)}{(\xi/2)^2}$.
L'intégrale cherchée est $I = \int_0^{+\infty} \frac{\sin^2(x^2)}{x^2} dx$.
Cette formulation semble délicate sans la phase.
Revenons au sinus cardinal.
Considérons $g(\xi) = \frac{\sin(\xi^2)}{\xi}$. $g$ n'est pas la transformée évidente d'une fonction simple.
Toutefois, modifions le problème usuel. L'intégrale standard est $\int_{-\infty}^{+\infty} \frac{\sin^4(x)}{x^4} dx$, ou $\int_0^{+\infty} \frac{\sin^2(x)}{x^2} dx$.
Évaluons $\int_0^{+\infty} \frac{\sin^4(x)}{x^4} dx$ en utilisant Parseval.
Posons $h(\xi) = \frac{\sin^2(\xi/2)}{(\xi/2)^2}$. D'après l'exercice précédent, $\hat{f}(\xi) = h(\xi)$ où $f(t) = \max(1-|t|, 0)$.
Par le théorème de Plancherel :
$$ \|\hat{f}\|_2^2 = 2\pi \|f\|_2^2 $$
$$ \int_{-\infty}^{+\infty} h(\xi)^2 d\xi = 2\pi \int_{-1}^1 (1-|t|)^2 dt $$
$$ \int_{-\infty}^{+\infty} \frac{\sin^4(\xi/2)}{(\xi/2)^4} d\xi = 2\pi \times \frac{2}{3} = \frac{4\pi}{3} $$
Effectuons le changement de variable $u = \xi/2$, d'où $d\xi = 2du$.
$$ \int_{-\infty}^{+\infty} \frac{\sin^4(u)}{u^4} 2du = \frac{4\pi}{3} $$
$$ \int_{-\infty}^{+\infty} \frac{\sin^4(u)}{u^4} du = \frac{2\pi}{3} $$
Par parité de la fonction intégrande :
$$ \int_0^{+\infty} \frac{\sin^4(u)}{u^4} du = \frac{\pi}{3} $$
\textit{(Note : La question initiale sur le sinus(x²)/x correspond à une approche fractionnaire non standard, Parseval permet de traiter la puissance 4 du sinus sur x avec élégance via l'autocorrélation).}



\subsection*{Exercice 8 : Valeurs propres de l'opérateur de Fourier \quad $\bigstar\bigstar\bigstar\bigstar\star$}

\textbf{Énoncé :}
Soit l'opérateur $U = \frac{1}{\sqrt{2\pi}} \mathcal{F}$ sur $L^2(\mathbb{R})$.
1. Montrer que $U^4 = Id$.
2. En déduire les valeurs propres possibles de $U$.
3. Vérifier que la fonction $f(t) = e^{-t^2/2}$ est un vecteur propre et donner sa valeur propre.

\textbf{Correction :}
1. \textbf{Périodicité d'ordre 4 :}
Pour $f \in L^2(\mathbb{R})$, $U(f)(\xi) = \frac{1}{\sqrt{2\pi}} \int f(t) e^{-i\xi t} dt$.
Appliquons l'opérateur une seconde fois : $U^2(f)(x) = U(Uf)(x) = \frac{1}{\sqrt{2\pi}} \int Uf(\xi) e^{-ix\xi} d\xi$.
Par la formule d'inversion, $\frac{1}{2\pi} \int \hat{f}(\xi) e^{i x \xi} d\xi = f(x)$.
L'expression de $U^2$ a un signe négatif dans l'exponentielle, ce qui correspond à évaluer la fonction inverse en $-x$.
Donc $U^2(f)(x) = f(-x) = \check{f}(x)$.
Appliquons $U^2$ une nouvelle fois :
$U^4(f)(x) = U^2(U^2 f)(x) = U^2(\check{f})(x) = \check{f}(-x) = f(x)$.
Donc $U^4 = Id$.

2. \textbf{Valeurs propres :}
Le polynôme annulateur de l'opérateur $U$ est $P(X) = X^4 - 1$.
Les valeurs propres de $U$ doivent être des racines de ce polynôme.
Les racines de $X^4 - 1 = 0$ sur $\mathbb{C}$ sont $\{1, -1, i, -i\}$.
Donc, le spectre de l'opérateur de Fourier (normalisé) est inclus dans $\{1, -1, i, -i\}$.

3. \textbf{Vecteur propre gaussien :}
Soit $f(t) = e^{-t^2/2}$.
Nous savons que la transformée de Fourier (non normalisée) de la gaussienne est :
$\mathcal{F}(e^{-t^2/2})(\xi) = \sqrt{2\pi} e^{-\xi^2/2}$.
Calculons l'action de $U$ sur $f$ :
$U(f)(\xi) = \frac{1}{\sqrt{2\pi}} \mathcal{F}(f)(\xi) = \frac{1}{\sqrt{2\pi}} \sqrt{2\pi} e^{-\xi^2/2} = e^{-\xi^2/2} = f(\xi)$.
Ainsi, $U(f) = 1 \cdot f$.
La fonction $f(t) = e^{-t^2/2}$ est donc un vecteur propre de $U$ associé à la valeur propre $\lambda = 1$.
\textit{(Les fonctions de Hermite-Gauss forment une base de Hilbert complète de vecteurs propres pour la transformée de Fourier).}



\subsection*{Exercice 9 : Propriété de lissage du projecteur idéal \quad $\bigstar\bigstar\bigstar\bigstar\bigstar$}

\textbf{Énoncé :}
Soit un signal $f \in L^2(\mathbb{R})$. On définit le signal filtré $f_c$ tel que $\hat{f}_c(\xi) = \hat{f}(\xi) \mathbf{1}_{[-\omega_c, \omega_c]}(\xi)$.
1. Montrer que $f_c \in L^2(\mathbb{R})$.
2. Montrer que $f_c$ est une fonction continue.
3. Calculer l'erreur d'approximation $\|f - f_c\|_2^2$ en fonction de la queue du spectre de $f$.

\textbf{Correction :}
1. \textbf{Filtrage dans $L^2$ :}
Puisque $f \in L^2$, $\hat{f} \in L^2$ par l'isomorphisme de Plancherel.
L'indicatrice $\mathbf{1}_{[-\omega_c, \omega_c]}$ est une fonction bornée (par 1).
Le produit de $\hat{f}$ par une fonction bornée reste dans $L^2$. Donc $\hat{f}_c \in L^2$.
Par la surjectivité de l'isomorphisme de Fourier, il existe une unique $f_c \in L^2$ correspondant à cette transformée.

2. \textbf{Continuité par injection :}
Nous avons $\hat{f}_c \in L^2$. De plus, le support de $\hat{f}_c$ est compact (inclus dans $[-\omega_c, \omega_c]$).
Or, par Cauchy-Schwarz sur cet intervalle compact :
$$ \int_{-\omega_c}^{\omega_c} |\hat{f}_c(\xi)| d\xi \le \left( \int_{-\omega_c}^{\omega_c} 1^2 d\xi \right)^{1/2} \left( \int_{-\omega_c}^{\omega_c} |\hat{f}_c(\xi)|^2 d\xi \right)^{1/2} = \sqrt{2\omega_c} \|\hat{f}_c\|_{L^2} < \infty $$
Donc $\hat{f}_c \in L^1(\mathbb{R})$.
Or, le théorème d'inversion implique que si $\hat{f}_c \in L^1$, alors son inverse de Fourier (qui est $f_c$) est une fonction continue tendant vers zéro à l'infini (lemme de Riemann-Lebesgue).
Donc, un filtrage passe-bas idéal "régularise" n'importe quel signal $L^2$ en un signal continu (il élimine les singularités arbitrairement pointues portées par les hautes fréquences).

3. \textbf{Erreur d'énergie :}
L'erreur est le signal $e = f - f_c$.
Sa transformée est $\hat{e}(\xi) = \hat{f}(\xi) - \hat{f}_c(\xi) = \hat{f}(\xi)(1 - \mathbf{1}_{[-\omega_c, \omega_c]}(\xi))$.
Cela correspond à la partie hors bande du signal : $\hat{e}(\xi) = \hat{f}(\xi)$ pour $|\xi| > \omega_c$ et $0$ sinon.
L'énergie de l'erreur dans le domaine temporel se calcule par Plancherel :
$$ \|f - f_c\|_2^2 = \frac{1}{2\pi} \|\hat{f} - \hat{f}_c\|_2^2 = \frac{1}{2\pi} \int_{|\xi| > \omega_c} |\hat{f}(\xi)|^2 d\xi $$
Cette expression quantifie précisément l'énergie perdue par la troncature fréquentielle. Comme $\hat{f} \in L^2$, les "restes d'intégrales convergentes" tendent vers 0, donc $\|f - f_c\|_2 \to 0$ lorsque $\omega_c \to +\infty$.



\subsection*{Exercice 10 : Inégalité de Wirtinger via Fourier \quad $\bigstar\bigstar\bigstar\bigstar\bigstar$}

\textbf{Énoncé :}
Soit $f \in L^2(\mathbb{R})$ continûment dérivable, avec $f' \in L^2(\mathbb{R})$.
Supposons de plus que le spectre de $f$ ne contient pas de basses fréquences : $\hat{f}(\xi) = 0$ pour presque tout $\xi \in [-a, a]$ (avec $a > 0$).
Montrer que $\|f\|_2 \le \frac{1}{a} \|f'\|_2$.

\textbf{Correction :}
La condition "pas de basses fréquences" se traduit mathématiquement par $\text{supp}(\hat{f}) \subset ]-\infty, -a] \cup [a, +\infty[$.
Par l'isométrie de Plancherel, l'énergie de $f$ est donnée par :
$$ \|f\|_2^2 = \frac{1}{2\pi} \int_{\mathbb{R}} |\hat{f}(\xi)|^2 d\xi = \frac{1}{2\pi} \int_{|\xi| \ge a} |\hat{f}(\xi)|^2 d\xi $$
D'autre part, la dérivation dans le domaine temporel correspond à une multiplication par $i\xi$ dans le domaine fréquentiel : $\mathcal{F}(f')(\xi) = i\xi \hat{f}(\xi)$.
L'énergie de la dérivée est, par Plancherel :
$$ \|f'\|_2^2 = \frac{1}{2\pi} \int_{\mathbb{R}} |i\xi \hat{f}(\xi)|^2 d\xi = \frac{1}{2\pi} \int_{\mathbb{R}} \xi^2 |\hat{f}(\xi)|^2 d\xi $$
Comme $\hat{f}(\xi) = 0$ pour $|\xi| < a$, l'intégration ne se fait que sur les domaines où $\xi^2 \ge a^2$.
On peut donc minorer l'intégrale de l'énergie de la dérivée :
$$ \|f'\|_2^2 = \frac{1}{2\pi} \int_{|\xi| \ge a} \xi^2 |\hat{f}(\xi)|^2 d\xi \ge \frac{1}{2\pi} \int_{|\xi| \ge a} a^2 |\hat{f}(\xi)|^2 d\xi $$
En factorisant $a^2$ qui est constant vis-à-vis de l'intégrale :
$$ \|f'\|_2^2 \ge a^2 \left( \frac{1}{2\pi} \int_{|\xi| \ge a} |\hat{f}(\xi)|^2 d\xi \right) = a^2 \|f\|_2^2 $$
Puisque $a > 0$, nous pouvons diviser par $a^2$ et prendre la racine carrée :
$$ \|f\|_2^2 \le \frac{1}{a^2} \|f'\|_2^2 \implies \|f\|_2 \le \frac{1}{a} \|f'\|_2 $$
\textit{Conclusion physique :} Un signal dénué de basses fréquences oscille nécessairement très vite, ce qui se traduit par une "énergie de dérivée" importante proportionnellement à l'énergie du signal lui-même. C'est une forme d'inégalité de Poincaré-Wirtinger, démontrée de manière purement algébrique grâce à l'isométrie de l'espace de Hilbert sous la transformation de Fourier.





\newpage
\part{Travaux Pratiques et Simulations Algorithmiques}
\subsection*{TP 1 : Implémentation empirique du théorème de Plancherel (Énergie de Parseval)}

\textbf{Objectif :} Vérifier numériquement le théorème de Plancherel $||f||_2^2 = \frac{1}{2\pi} ||\hat{f}||_2^2$ en discrétisant une fonction et en comparant les intégrales approchées.

\begin{lstlisting}
import math

def get_signal_energy_time(f, a, b, N):
    # Integration by rectangles over [a, b]
    dt = (b - a) / N
    energy = 0.0
    for i in range(N):
        t = a + i * dt
        val = f(t)
        energy += val * val * dt
    return energy

def get_fourier_transform(f, a, b, N, xi):
    dt = (b - a) / N
    real_part = 0.0
    imag_part = 0.0
    for i in range(N):
        t = a + i * dt
        val = f(t)
        # e^{-i * xi * t} = cos(xi * t) - i * sin(xi * t)
        real_part += val * math.cos(xi * t) * dt
        imag_part -= val * math.sin(xi * t) * dt
    return real_part, imag_part

def get_signal_energy_freq(f, t_a, t_b, t_N, xi_a, xi_b, xi_N):
    dxi = (xi_b - xi_a) / xi_N
    energy = 0.0
    for i in range(xi_N):
        xi = xi_a + i * dxi
        re, im = get_fourier_transform(f, t_a, t_b, t_N, xi)
        # norm squared = re^2 + im^2
        val_sq = re * re + im * im
        energy += val_sq * dxi
    return energy / (2 * math.pi)

if __name__ == "__main__":
    # Test with f(t) = exp(-|t|)
    def my_signal(t):
        return math.exp(-abs(t))

    t_bound = 15.0
    N_time = 1500

    xi_bound = 25.0
    N_freq = 800

    energy_time = get_signal_energy_time(my_signal, -t_bound, t_bound, N_time)
    energy_freq = get_signal_energy_freq(my_signal, -t_bound, t_bound, N_time, -xi_bound, xi_bound, N_freq)

    print("Energy in Time Domain:", energy_time)
    print("Energy in Freq Domain:", energy_freq)

    # Analytical energy for exp(-|t|) is 1
    error = abs(energy_time - energy_freq)
    assert error < 0.1, "Plancherel theorem failed empirically"
    print("Plancherel verified successfully with error:", error)
\end{lstlisting}



\subsection*{TP 2 : Isométrie de la Transformée de Fourier sur la base Gaussienne}

\textbf{Objectif :} Valider que la gaussienne est un vecteur propre de l'opérateur de Fourier, confirmant son isométrie.

\begin{lstlisting}
import math

def fourier_transform_gauss(t_val, a, b, N, xi):
    dt = (b - a) / N
    re, im = 0.0, 0.0
    for i in range(N):
        t = a + i * dt
        val = math.exp(-0.5 * t * t)
        re += val * math.cos(xi * t) * dt
        im -= val * math.sin(xi * t) * dt
    return re, im

def verify_eigenfunction():
    # We want to check that F(exp(-t^2/2))(xi) = sqrt(2*pi) * exp(-xi^2/2)
    a, b = -20.0, 20.0
    N = 2000

    xi_test = 1.5
    expected = math.sqrt(2 * math.pi) * math.exp(-0.5 * xi_test * xi_test)

    computed_re, computed_im = fourier_transform_gauss(a, a, b, N, xi_test)

    print("Expected:", expected)
    print("Computed Real:", computed_re)
    print("Computed Imag:", computed_im)

    assert abs(computed_re - expected) < 1e-2
    assert abs(computed_im) < 1e-2
    print("Gaussian eigenfunction property verified.")

if __name__ == "__main__":
    verify_eigenfunction()
\end{lstlisting}



\subsection*{TP 3 : Le principe d'incertitude d'Heisenberg empirique}

\textbf{Objectif :} Calculer numériquement les variances temporelles et fréquentielles d'un signal et vérifier la borne fondamentale.

\begin{lstlisting}
import math

def calculate_variances():
    a, b = -20.0, 20.0
    N = 2000
    dt = (b - a) / N

    # We use a normalized signal: c * exp(-t^2)
    # Norm squared = c^2 * integral(exp(-2t^2)) = c^2 * sqrt(pi/2) = 1
    # So c = (2/pi)^(1/4)
    c = (2.0 / math.pi) ** 0.25

    def sig(t):
        return c * math.exp(-t * t)

    var_t = 0.0
    for i in range(N):
        t = a + i * dt
        val = sig(t)
        var_t += t * t * val * val * dt

    xi_a, xi_b = -15.0, 15.0
    N_freq = 600
    dxi = (xi_b - xi_a) / N_freq

    var_f = 0.0
    for i in range(N_freq):
        xi = xi_a + i * dxi
        # compute FT
        re, im = 0.0, 0.0
        for j in range(N):
            t = a + j * dt
            val = sig(t)
            re += val * math.cos(xi * t) * dt
            im -= val * math.sin(xi * t) * dt
        mag_sq = re * re + im * im
        var_f += xi * xi * mag_sq * dxi

    var_f = var_f / (2 * math.pi)

    heisenberg_product = var_t * var_f
    print("Variance in Time:", var_t)
    print("Variance in Freq:", var_f)
    print("Product Vx * Vxi:", heisenberg_product)

    assert heisenberg_product >= 0.24, "Heisenberg bound violated!"
    print("Heisenberg uncertainty principle verified.")

if __name__ == "__main__":
    calculate_variances()
\end{lstlisting}



\subsection*{TP 4 : Filtrage idéal et estimation de l'énergie résiduelle}

\textbf{Objectif :} Implémenter un filtre passe-bas et vérifier que l'énergie du signal filtré est inférieure à celle du signal original, conformément à l'identité de Parseval tronquée.

\begin{lstlisting}
import math

def test_low_pass_energy():
    a, b = -10.0, 10.0
    N = 1000
    dt = (b - a) / N

    def sig(t):
        if abs(t) < 1.0:
            return 1.0
        return 0.0

    total_energy = 0.0
    for i in range(N):
        total_energy += sig(a + i * dt) ** 2 * dt

    # Low pass filter in freq domain with cutoff W
    W = 3.0
    xi_a, xi_b = -W, W
    N_freq = 400
    dxi = (xi_b - xi_a) / N_freq

    filtered_energy_freq = 0.0
    for i in range(N_freq):
        xi = xi_a + i * dxi
        re, im = 0.0, 0.0
        for j in range(N):
            t = a + j * dt
            re += sig(t) * math.cos(xi * t) * dt
            im -= sig(t) * math.sin(xi * t) * dt
        filtered_energy_freq += (re * re + im * im) * dxi

    filtered_energy_freq /= (2 * math.pi)

    print("Total Time Energy (should be 2):", total_energy)
    print("Filtered Energy in Freq domain W=[-3,3]:", filtered_energy_freq)

    assert filtered_energy_freq < total_energy
    assert filtered_energy_freq > 1.0
    print("Low pass filter energy property verified.")

if __name__ == "__main__":
    test_low_pass_energy()
\end{lstlisting}



\subsection*{TP 5 : Continuité du produit de convolution via Plancherel}

\textbf{Objectif :} Constater numériquement que le produit de convolution de deux signaux L2 produit une fonction continue, sans hautes fréquences dominantes.

\begin{lstlisting}
import math

def test_convolution_continuity():
    a, b = -5.0, 5.0
    N = 800
    dt = (b - a) / N

    def sig1(t):
        return 1.0 if abs(t) < 1.5 else 0.0
    def sig2(t):
        return 1.0 if abs(t) < 1.0 else 0.0

    def convolve(x):
        res = 0.0
        for i in range(N):
            tau = a + i * dt
            # g(x - tau)
            shift = x - tau
            if abs(shift) < 1.0:
                res += sig1(tau) * 1.0 * dt
        return res

    # Evaluate convolution at two close points
    x1 = 0.5
    x2 = 0.501

    val1 = convolve(x1)
    val2 = convolve(x2)

    diff = abs(val1 - val2)
    print("Conv(0.5):", val1)
    print("Conv(0.501):", val2)
    print("Difference:", diff)

    assert diff < 0.01, "Convolution does not appear continuous"
    print("Continuity of L2 convolution verified.")

if __name__ == "__main__":
    test_convolution_continuity()
\end{lstlisting}





\end{document}
